\documentclass{article}
\usepackage{graphicx}
\usepackage{booktabs}

\title{Antigravity Findings v2.0: A Bayesian Analysis of Temporal Leakage in LLMs}
\author{The Antigravity Community}
\date{July 2026}

\begin{document}
\maketitle

\begin{abstract}
We present a consolidated analysis of independent external reproductions (N=25) of the ConstraintBench-100 dataset. 
Using Bayesian inference, we model the true temporal leakage rates of major frontier models, confirming that while larger models 
exhibit lower baseline leakage, the phenomenon remains statistically significant across all evaluated architectures.
\end{abstract}

\section{Results}
Table 1 presents the 95\% Credible Intervals for temporal leakage rates across models.

\begin{table}[h]
\centering
\begin{tabular}{lccc}
\toprule
Model & $N_{evals}$ & Mean Leakage (\%) & 95\% CI (\%) \\
\midrule
llama3.1:8b-instruct-q4\_K\_M & 4 & 33.9 & [28.2, 39.5] \\
gemma2:9b-instruct & 6 & 27.8 & [22.3, 33.3] \\
llama3:70b-instruct & 5 & 15.4 & [9.0, 21.8] \\
gpt-4o-2024-05-13 & 6 & 10.8 & [4.1, 17.5] \\
claude-3.5-sonnet & 4 & 8.1 & [5.8, 10.5] \\
\bottomrule
\end{tabular}
\caption{Bayesian Consolidation of Community Reproductions}
\end{table}

\section{Conclusion}
The temporal leakage phenomenon is highly reproducible. Future work must investigate structural interventions to the attention mechanism.

\end{document}
